\section{What sparsity makes available}
\label{sec:measurement}
We consider full-sequence causal inference at sequence length $T$. For each
matrix-product operation $o$, let $N_o$ count scalar multiplications and let
$Z_o$ count those for which at least one operand is exactly zero. We count a
multiplication once when both operands are zero. The transformer-block
operations are the query/key/value projections, attention output projection,
two feed-forward projections, and the attention products $QK^\top$ and $PV$.
The vocabulary projection contributes its full dense count, $N_{\rm head}$.
We define the \emph{model-wide activation sparsity opportunity} as
\begin{equation}
 \Rmodel =
 \frac{\sum_{o\in{\rm blocks}} Z_o}
      {\sum_{o\in{\rm blocks}} N_o+N_{\rm head}},
 \qquad
 \Rblock =
 \frac{\sum_{o\in{\rm blocks}} Z_o}{\sum_{o\in{\rm blocks}} N_o}.
 \label{eq:opportunity}
\end{equation}
The reported interventions act on activations; the counter uses the actual
operands, including any naturally occurring zeros. Integer counts are pooled
over all layers, batches, and validation sequences before division.

This definition connects local zeros to the operations that use them.
For example, a zero feed-forward hidden activation affects the following
projection. A zero key component can participate in multiple causal
query--key pairs. For attention products, we count the union of zero-operand
events over valid causal positions using the post-rotation queries and keys.
Appendix~\ref{app:counting} gives the counts and denominator explicitly.
Layer normalization, softmax, elementwise transformations, data movement, and
kernel-launch costs fall outside this multiplication accounting.

Cross-architecture comparisons additionally depend on the block share
$s_{\rm block}$ of the counted work:
\begin{equation}
  \Rmodel = s_{\rm block}\Rblock,\qquad
  s_{\rm block}=\frac{\sum_{o\in{\rm blocks}}N_o}
  {\sum_{o\in{\rm blocks}}N_o+N_{\rm head}}.
  \label{eq:block-share}
\end{equation}
This identity helps separate learned block sparsity from architectural
differences. Runtime measurements then test how much of the opportunity a
particular implementation realizes.
